\documentclass{article}

\usepackage[final]{neurips_2020}

\usepackage[utf8]{inputenc}
\usepackage[T1]{fontenc}
\usepackage{hyperref}
\usepackage{url}
\usepackage{booktabs}
\usepackage{amsfonts}
\usepackage{amsmath}
\usepackage{nicefrac}
\usepackage{microtype}
\usepackage{graphicx}
\usepackage{xcolor}
\usepackage{multirow}

% \note{...} prints TODO reminders in blue. Delete notes before final submission.
\newcommand{\note}[1]{\textcolor{blue}{{#1}}}

\title{
  Difficulty-Aware Language Routing for Korean Mathematical Reasoning \\
  \vspace{1em}
  \small{\normalfont Korea University COSE461 Final Project}
}

\author{
  \note{Your Name} \\
   Department of Computer Science \\
   \note{Team N} \\
   \note{2022320317}
   \And
   \note{Teammate Name} \\
   Department of Computer Science \\
   \note{Team N} \\
   \note{20XXXXXXXX}
}

\begin{document}

\maketitle

\begin{abstract}
Small language models reason markedly better in English than in Korean
on mathematical word problems. We trace part of this gap to the
distillation pipeline itself: even a strong teacher produces
lower-quality chain-of-thought (CoT) in Korean (85.7\% correct) than in
English (97.5\%), so naively distilling Korean CoT caps the student at a
degraded signal. We ask whether a model's stronger English reasoning can
be leveraged to improve its Korean reasoning. We introduce
\textbf{Difficulty-Aware Language Routing (DALR)}, a data-construction
method that decides, per problem, whether to learn from a Korean or an
English teacher CoT based on teacher correctness: English CoT is used
only as a ``bridge'' on problems the Korean teacher fails. A controlled
ablation shows that DALR's improvement comes from \emph{where} English
reasoning is routed (hard problems) rather than from adding more data.
We further combine specialist students via model soup
\citep{wortsman2022soup} and a cross-lingual ensemble at inference. On
HRM8K (Korean) and GSM8K (English) with Qwen2.5-3B, DALR achieves the
best single-model Korean accuracy, model soup resolves the
Korean--English trade-off, and the cross-lingual ensemble yields the
best overall results.
\end{abstract}

% =====================================================================
\section{Introduction}
\label{sec:intro}

Large language models perform strong multi-step mathematical reasoning,
but small models (e.g.\ 3B parameters) lag far behind, and the gap is
worse outside English. For a 3B model (Qwen2.5-3B-Instruct), we observe a
large discrepancy between English and Korean: it solves \textbf{67.4\%}
of English problems (GSM8K) but only \textbf{52.4\%} of Korean problems
(HRM8K). This matters because most practical deployments cannot afford
large models, yet must serve non-English users.

A standard remedy is knowledge distillation: a strong teacher generates
chain-of-thought (CoT) solutions, and a small student is fine-tuned on
them. However, this strategy faces a subtle obstacle in the
cross-lingual setting: \emph{the teacher is itself weaker in Korean}.
In our data, the teacher produces correct CoT for \textbf{97.5\%} of
English problems but only \textbf{85.7\%} of Korean problems. Distilling
Korean CoT therefore transfers a noisier signal, while English continues
to improve, widening rather than closing the gap.

This motivates our central question: \emph{can a model's (and a
teacher's) stronger English reasoning be used to improve Korean
mathematical reasoning?} Our working assumption is that mathematical
\emph{logic} is largely language-agnostic---the solution steps are the
same whether a problem is posed in Korean or English---so connecting
English reasoning patterns to Korean problems should help.

We pursue this through \emph{quality-aware aggregation} of cross-lingual
signals at three levels. At the \textbf{data level}, our main
contribution \textbf{DALR} routes each problem to a Korean or English
teacher CoT based on teacher correctness. At the \textbf{model level}, we
average the weights of specialist students (model soup,
\citealp{wortsman2022soup}). At the \textbf{inference level}, we
majority-vote across models that reason in different languages.

\paragraph{Contributions.}
\begin{itemize}
  \item We propose \textbf{DALR}, a difficulty-aware data-construction
  method for cross-lingual CoT distillation, and show via a controlled
  ablation (random vs.\ routed English bridges) that its gains stem from
  \emph{routing}, not data scaling.
  \item We apply and analyze model soup and cross-lingual ensembling in
  the Korean math setting, identifying a Korean--English trade-off
  induced by DALR and showing how aggregation resolves it.
  \item We provide a statistical analysis (bootstrap confidence
  intervals, McNemar tests) and a per-difficulty breakdown clarifying
  where each method helps.
\end{itemize}

% =====================================================================
\section{Related Work}
\label{sec:related}

\paragraph{Chain-of-thought and its distillation.}
Chain-of-thought prompting elicits step-by-step reasoning
\citep{wei2022cot}, and a line of work distills such reasoning into
smaller models. \note{Add 1--2 CoT-distillation citations.}

\paragraph{Cross-lingual reasoning.}
Prior work transfers reasoning across languages via translation or
cross-lingual instruction tuning. \note{Cite e.g.\ xCoT / MGSM.} Unlike
methods that augment prompts at inference, DALR selects the training CoT
\emph{language} per problem according to teacher quality.

\paragraph{Weight averaging and self-consistency.}
Model soup averages the weights of multiple fine-tuned models to obtain a
single model at no extra inference cost \citep{wortsman2022soup}.
Self-consistency samples multiple chains from one model and votes
\citep{wang2022selfconsistency}. We instead ensemble \emph{distinct}
models trained in different languages, exploiting their differing error
patterns.

% =====================================================================
\section{Approach}
\label{sec:approach}

\subsection{Overview}
We train a small student with several recipes that differ only in the
language composition of their distillation data, producing complementary
``specialists.'' We then aggregate cross-lingual signal at three levels:
data (DALR), model weights (soup), and inference (ensemble).
\note{Add Figure 1: the three-level pipeline diagram.}

\subsection{Difficulty-Aware Language Routing (DALR)}
\label{sec:dalr}
Let a training problem have a Korean teacher CoT and an English teacher
CoT, each labeled correct or incorrect by answer matching. For each
Korean problem we route:
\begin{itemize}
  \item Korean correct $\rightarrow$ train on the \textbf{Korean CoT}
  (natural Korean reasoning);
  \item Korean incorrect, English correct $\rightarrow$ train on the
  \textbf{English CoT} as a \emph{bridge} (Korean question, English
  solution);
  \item both incorrect $\rightarrow$ \textbf{discard} (likely noise).
\end{itemize}
Teacher failure in Korean is used as a signal that the problem is hard to
reason about in Korean; a reliable English solution is then a better
target. This yields 6{,}407 Korean and 920 English-bridge examples
(7{,}327 total), which we call setup~F.

\subsection{Model Soup}
\label{sec:soup}
Our recipes yield specialists with different strengths (English-strong,
Korean-strong, balanced). We combine them by element-wise averaging of
their LoRA weights \citep{wortsman2022soup}, producing a single model at
$1\times$ inference cost.

\subsection{Cross-Lingual Ensemble}
\label{sec:clsc}
At inference we run several models on the same problem and take a
majority vote over final answers, breaking ties by a designated model.
Because the models reason in different languages, their errors are
decorrelated, improving ensemble accuracy.

\subsection{Baselines}
We compare against: A (base, no fine-tuning), B (English-only CoT), C
(Korean-only CoT), D (bilingual mix), and E (two-stage English-then-
Korean). F is DALR; F\_random is the DALR ablation that places the same
920 English bridges on randomly chosen \emph{easy} (Korean-solved)
problems instead of hard ones.

% =====================================================================
\section{Experiments}
\label{sec:exp}

\subsection{Data}
Training problems come from GSM8K (7{,}473 grade-school math word
problems) and its Korean machine translation GSM8K-ko. The teacher
(Gemini) generates CoT in each language; we keep only answer-correct CoT
(English: 7{,}283; Korean: 6{,}407). We evaluate on \textbf{HRM8K}
(human-reviewed Korean translation of the GSM8K test set, 1{,}319
problems) and \textbf{GSM8K} (English, 1{,}319). HRM8K is our main
benchmark because its translation quality is higher than GSM8K-ko.

\subsection{Evaluation method}
We report exact-match accuracy on the final numeric answer, extracted by
a regular expression and compared to the gold answer with a small
numeric tolerance. We use single greedy decoding (no sampling) with a
zero-shot CoT system prompt and a repetition penalty of 1.1.

\subsection{Experimental details}
The student is Qwen2.5-3B-Instruct, fine-tuned with LoRA ($r{=}32$,
$\alpha{=}64$) for 3 epochs at learning rate $2\times10^{-4}$
\note{(E uses two stages; see appendix)}. We train on a single
RTX~5060~Ti. Model soups are formed by averaging LoRA adapters across
setups.

\subsection{Results}
\label{sec:results}

\begin{table}[t]
\centering
\small
\begin{tabular}{lcc}
\toprule
\textbf{Setup} & \textbf{HRM8K (KO)} & \textbf{GSM8K (EN)} \\
\midrule
A (Base)              & 53.9 & \note{TBD} \\
B (English CoT)       & 54.4 & \note{TBD} \\
C (Korean CoT)        & 59.5 & 70.7 \\
D (Bilingual mix)     & \note{TBD} & \note{TBD} \\
E (Two-stage)         & \note{TBD} & \note{TBD} \\
F (DALR)              & \textbf{61.8} & 69.4 \\
F\_random (ablation)  & 58.8 & 72.7 \\
\midrule
Model soup (B+C+D+E+F) & 60.1 & \textbf{78.6} \\
\midrule
Ensemble, 6 models      & 68.6$^\dagger$ & 81.6$^\dagger$ \\
Ensemble, 6 + soup      & \textbf{70.0}$^\dagger$ & \textbf{84.4}$^\dagger$ \\
\bottomrule
\end{tabular}
\caption{Accuracy (\%) on the full 1{,}319-problem test sets.
$^\dagger$ ensemble numbers are preliminary (500-sample); the full-set
update is pending. \note{Fill TBD once eval finishes.}}
\label{tab:main}
\end{table}

Table~\ref{tab:main} summarizes our results.

\paragraph{DALR improves Korean, and routing (not data) is the cause.}
F and F\_random use the \emph{same} 920 English bridges and differ only
in placement. On HRM8K, F (61.8) exceeds F\_random (58.8) by 3.0 points
(McNemar $p{=}0.030$), and F\_random even drops below the Korean-only
baseline C (59.5). Adding English data \emph{at random} hurts; routing it
to hard problems helps. This isolates routing as the source of the gain.

\paragraph{A Korean--English trade-off, resolved by aggregation.}
DALR raises Korean but slightly lowers English (F: 69.4 vs.\ C: 70.7;
$p{=}0.016$ vs.\ F\_random). Model soup recovers English strongly (78.6,
$p{<}0.001$ vs.\ F) while keeping competitive Korean, giving the best
single deployable model. The cross-lingual ensemble attains the best
overall accuracy in both languages.

% =====================================================================
\section{Analysis}
\label{sec:analysis}

\paragraph{Where does DALR help?}
Stratifying HRM8K problems by how many of our models solve them, F is
strongest on \emph{medium-difficulty} problems
\note{(insert per-difficulty numbers from final\_report.py)}, consistent
with the intuition that English bridges help where Korean-only reasoning
is unreliable, while contributing little on problems that are either
trivial or unsolvable by all models.

\paragraph{Statistical robustness.}
We report bootstrap 95\% confidence intervals for all setups and McNemar
tests for key comparisons \note{(move final numbers here after 1,319
eval)}. The DALR-vs-Korean-only gap on HRM8K is positive but modest; the
routing ablation (F vs.\ F\_random) and the soup's English gains are the
statistically strongest effects.

% =====================================================================
\section{Conclusion}
\label{sec:conclusion}
We studied how to improve Korean mathematical reasoning in a small model
by leveraging stronger English reasoning. Our main contribution, DALR,
routes teacher CoT language per problem by difficulty; a controlled
ablation shows the benefit comes from routing rather than added data.
Combined with model soup and cross-lingual ensembling, the approach
improves both Korean and English accuracy over baselines.
\textbf{Limitations.} HRM8K is a translation of GSM8K rather than
natively Korean math, so gains may partly reflect translated-problem
solving; the DALR-vs-Korean-only margin is modest; and we use a single
teacher and a single 3B student.

\bibliographystyle{unsrt}
\bibliography{references}

\newpage
\appendix

\section{Appendix: Team contributions}
\note{1--2 sentences per member.}

\end{document}
